\section{Évolution de l'émetteur : L'itération OpenHD}

L'analyse des performances de l'architecture initiale basée sur la Nvidia Jetson Nano a mis en exergue un 
goulot d'étranglement matériel majeur : l'Image Signal Processor (ISP)\index{ISP} de la caméra, qui impose 
une file d'attente logicielle inévitable générant environ 130 ms de latence résiduelle (à 30 FPS). Afin de 
s'affranchir de cette limitation et de repousser les limites du temps réel, une évolution complète du 
système embarqué a été étudiée.

\subsection{Changement de paradigme matériel : La Raspberry Pi 4B}
Pour reprendre le contrôle total sur le flux vidéo dès la sortie du capteur CMOS\index{Capteur CMOS}, la 
Jetson Nano a été remplacée par une \textbf{Raspberry Pi 4B}\index{Raspberry Pi} agissant comme nouvel 
émetteur. 

Bien que la Jetson soit théoriquement plus puissante pour l'encodage H.265, l'écosystème logiciel de la 
Raspberry Pi offre un avantage critique pour la vidéo basse latence. Sous Linux standard, l'accès à la 
caméra passe par les API modernes telles que \texttt{libcamera} ou \texttt{V4L2}\index{V4L2}, qui implémentent 
des tampons de mémoire (\textit{buffers}) complexes pour garantir la qualité de l'image (Exposition, Balance 
des blancs) au détriment de la vitesse.

\subsection{L'écosystème OpenHD et l'API MMAL}
Au lieu de redévelopper une couche logicielle d'acquisition complète, le choix s'est porté sur l'utilisation 
du système d'exploitation du projet \textit{OpenHD}\index{OpenHD}. OpenHD est un projet \textit{open-source} 
de référence dans le domaine de la transmission vidéo numérique FPV.

Bien que ce projet ne traite pas directement de l'intégration avec notre modem Wi-Fi HaLow via SPI, son OS 
est optimisé pour un accès \textit{Bare-Metal} aux périphériques. L'atout majeur de l'OS OpenHD réside dans 
son utilisation de l'API \textbf{MMAL} (\textit{Multimedia Abstraction Layer}). 
Bien que considérée comme dépréciée par la fondation Raspberry Pi au profit de \texttt{libcamera}, l'API MMAL 
permet un accès direct (\textit{Zero-Copy}) au capteur de la caméra et à l'encodeur matériel (H.264) du 
processeur Broadcom, contournant intégralement les tampons de l'ISP.

\subsection{Intégration logicielle (SPI et Capture)}
Le portage du système a nécessité la réécriture du programme de transfert C. 
Au lieu de dépendre d'un pipeline GStreamer\index{GStreamer} externe générant un délai, le nouveau programme 
d'émission s'interface directement avec l'encodeur MMAL. 

La séquence de fonctionnement adaptée est la suivante :
\begin{enumerate}
    \item \textbf{Activation de l'interface matérielle :} Chargement du module noyau SPI (\texttt{sudo modprobe spidev}).
    \item \textbf{Configuration du capteur :} Forçage matériel de la caméra en mode basse latence 
    afin d'annuler les traitements d'image post-capture : 
    \begin{lstlisting}[language=bash]
v4l2-ctl -d /dev/video0 --set-ctrl low_latency_mode=1
    \end{lstlisting}
    \item \textbf{Synchronisation MCU :} Le mécanisme de \textit{Handshake} matériel (GPIO 78 vers 
    GPIO 15 du STM32\index{STM32}) développé pour la Jetson a été conservé et mappé sur le \textit{header} 
    GPIO de la Raspberry Pi pour éviter le débordement de la mémoire du modem.
\end{enumerate}

\subsection{Objectifs et perspectives de l'itération}
Cette itération architecturale déplace l'effort de la puissance de compression (Jetson H.265) vers la 
maîtrise pure de la chaîne d'acquisition (Raspberry Pi MMAL). L'objectif théorique de cette modification 
est l'oblitération des 130 ms de latence introduites par l'ISP de la caméra. Associée à notre couche 
réseau Wi-Fi HaLow optimisée (UDP Broadcast, MCS 2, No-ACK), cette configuration doit permettre 
d'atteindre une latence \textit{Glass-to-Glass} inférieure à 60 ms, rivalisant ainsi directement 
avec les systèmes numériques commerciaux opérant sur 5.8 GHz, tout en offrant une portée et une 
pénétration d'obstacles drastiquement supérieures.